%%%%%%%%%%%%%%%%%%%%%%%%%%%%%%%%%%%%%%%
% Wenneker Resume/CV
% LaTeX Template
% Version 1.1 (19/6/2016)
%
% This template has been downloaded from:
% http://www.LaTeXTemplates.com
%
% Original author:
% Frits Wenneker (http://www.howtotex.com) with extensive modifications by 
% Vel (vel@LaTeXTemplates.com)
%
% License:
% CC BY-NC-SA 3.0 (http://creativecommons.org/licenses/by-nc-sa/3.0/
%
%%%%%%%%%%%%%%%%%%%%%%%%%%%%%%%%%%%%%%

%----------------------------------------------------------------------------------------
%	PACKAGES AND OTHER DOCUMENT CONFIGURATIONS
%----------------------------------------------------------------------------------------

\documentclass[a4paper,12pt]{memoir} % Font and paper size

%%%%%%%%%%%%%%%%%%%%%%%%%%%%%%%%%%%%%%%%%
% Wenneker Resume/CV
% Structure Specification File
% Version 1.1 (19/6/2016)
%
% This file has been downloaded from:
% http://www.LaTeXTemplates.com
%
% Original author:
% Frits Wenneker (http://www.howtotex.com) with extensive modifications by 
% Vel (vel@latextemplates.com)
%
% License:
% CC BY-NC-SA 3.0 (http://creativecommons.org/licenses/by-nc-sa/3.0/)
%
%%%%%%%%%%%%%%%%%%%%%%%%%%%%%%%%%%%%%%%%%

%----------------------------------------------------------------------------------------
%	PACKAGES AND OTHER DOCUMENT CONFIGURATIONS
%----------------------------------------------------------------------------------------

\usepackage{XCharter} % Use the Bitstream Charter font
\usepackage[utf8]{inputenc} % Required for inputting international characters
\usepackage[T1]{fontenc} % Output font encoding for international characters

\usepackage[top=1cm,left=1cm,right=1cm,bottom=1cm]{geometry} % Modify margins

\usepackage{graphicx} % Required for figures

\usepackage{flowfram} % Required for the multi-column layout

\usepackage{url} % URLs

\usepackage[usenames,dvipsnames]{xcolor} % Required for custom colours

\usepackage{tikz} % Required for the horizontal rule

\usepackage[danish]{babel}

\usepackage{enumitem} % Required for modifying lists
\setlist{noitemsep,nolistsep} % Remove spacing within and around lists

\setlength{\columnsep}{\baselineskip} % Set the spacing between columns

% Define the left frame (sidebar)
\newflowframe{0.2\textwidth}{\textheight}{0pt}{0pt}[left]
\newlength{\LeftMainSep}
\setlength{\LeftMainSep}{0.2\textwidth}
\addtolength{\LeftMainSep}{1\columnsep}
 
% Small static frame for the vertical line
\newstaticframe{1.5pt}{\textheight}{\LeftMainSep}{0pt}
 
% Content of the static frame with the vertical line
\begin{staticcontents}{1}
\hfill
\tikz{\draw[loosely dotted,color=RoyalBlue,line width=1.5pt,yshift=0](0,0) -- (0,\textheight);}
\hfill\mbox{}
\end{staticcontents}
 
% Define the right frame (main body)
\addtolength{\LeftMainSep}{1.5pt}
\addtolength{\LeftMainSep}{1\columnsep}
\newflowframe{0.7\textwidth}{\textheight}{\LeftMainSep}{0pt}[main01]

\pagestyle{empty} % Disable all page numbering

\setlength{\parindent}{0pt} % Stop paragraph indentation

%----------------------------------------------------------------------------------------
%	NEW COMMANDS
%----------------------------------------------------------------------------------------

\newcommand{\userinformation}[1]{\renewcommand{\userinformation}{#1}} % Define a new command for the CV user's information that goes into the left column

\newcommand{\cvheading}[1]{{\Huge\bfseries\color{RoyalBlue} #1} \par\vspace{.6\baselineskip}} % New command for the CV heading
\newcommand{\cvsubheading}[1]{{\Large\bfseries #1} \bigbreak} % New command for the CV subheading

\newcommand{\Sep}{\vspace{1em}} % New command for the spacing between headings
\newcommand{\SmallSep}{\vspace{0.5em}} % New command for the spacing within headings

\newcommand{\aboutme}[2]{ % New command for the about me section
\textbf{\color{RoyalBlue} #1}~~#2\par\Sep
}
	
\newcommand{\CVSection}[1]{ % New command for the headings within sections
{\Large\textbf{#1}}\par
\SmallSep % Used for spacing
}

\newcommand{\CVItem}[2]{ % New command for the item descriptions
\textbf{\color{RoyalBlue} #1}\par
#2
\SmallSep % Used for spacing
}

\newcommand{\bluebullet}{\textcolor{RoyalBlue}{$\circ$}~~} % New command for the blue bullets
 % Include the file specifying document layout and packages

%----------------------------------------------------------------------------------------
%	NAME AND CONTACT INFORMATION 
%----------------------------------------------------------------------------------------

\userinformation{ % Set the content that goes into the sidebar of each page
\begin{flushright}
% Comment out this figure block if you don't want a photo
\includegraphics[width=0.6\columnwidth]{nathan.png}\\[\baselineskip] % Your photo
\small % Smaller font size
Nathan Hugh Barr \\ % Your name
\Sep
\url{nathanhbarr@gmail.com} \\ % Your email address
\Sep

tlf: +45 61 73 55 92 \\ % Your phone number
\Sep % Some whitespace
\textbf{Address} \\
Lysalleen 60 \\ % Address 1
4000 Roskilde \\ % Address 2
% \\ % Address 3
\Sep
\url{www.linkedin.com/in/nathanhbarr/} \\ % Your URL
\vfill % Whitespace under this block to push it up under the photo
\end{flushright}
}

%----------------------------------------------------------------------------------------

\begin{document}

\userinformation % Print your information in the left column

\framebreak % End of the first column

%----------------------------------------------------------------------------------------
%	HEADING
%----------------------------------------------------------------------------------------

\cvheading{Nathan Hugh Barr} % Large heading - your name

\cvsubheading{} % Subheading - your occupation/specialization

%----------------------------------------------------------------------------------------
%	ABOUT ME
%----------------------------------------------------------------------------------------

\aboutme{Profile}{My past jobs have been more focused on academia and teaching but i am a quick learner and i can immerse myself quickly in different projects. I am interested in data science and working with big databases as i enjoy exploring data sets and extracting knowledge from the data set. I am not shy when receiving constructive criticism and i am able to incorporate them immediately. I work well in a team environment.}


%----------------------------------------------------------------------------------------
%	EDUCATION
%----------------------------------------------------------------------------------------

\CVSection{Education}


%------------------------------------------------

\CVItem{2016 - 2019, Roskilde University}{Master of Science in Physics and Mathematics\\
\small Master Thesis: 
\begin{itemize}
\item Title: \textit{Optimising the thickness determination of homopolymers and diblock copolymers using optical spectral reflectance during solvent vapour annealing.}
\item In this master thesis, the homopolymers Polystyrene and Polyisoprene thin film wafers and the diblock-copolymer Polystyrene-b-polyisoprene thin film wafers thickness was investigated during solvent vapour annealing using toluene. The Fresnel models were fitted to the reflectance data measured using the optical spectral reflectance measurement technique and the fitting results and fitted thickness were analysed.   
\item Grade: 10
\end{itemize}}

%------------------------------------------------

\CVItem{2013 - 2016, Roskilde Universitet}{Bachelor of Natural Science in Physics and Mathematics\\
\small Bachelor project: 
\begin{itemize}
\item Title: \textit{Mathematical competencies gained in Danish High School Education from 1958 to the present - using integrals as a case study.}
\item In this project we attempted to measure which mathematical competences were strengthened and weakened through out three periods during a period from 1958 to the present. The used the formal theory of integrals as a case study. 
\item Grade: 10
\end{itemize}
}

%------------------------------------------------

\Sep % Extra whitespace after the end of a major section

%----------------------------------------------------------------------------------------
%	EXPERIENCE
%----------------------------------------------------------------------------------------

\CVSection{Experience and Volunteering}

%------------------------------------------------

\CVItem{January 2015 - December 2018, \textit{Private tutor}, My academy}{Mathematics and Physics tutor working primarily with students who found the subjects difficult. My job was to asses and correct their way of thinking in the subjects while creating a learning environment best suited to their learning outcomes. I am able to convey ideas clearly and coach the student in ways to handle problems and exam situations. I also have experience assessing the students pitfalls and planning new sessions.}


%------------------------------------------------

%\CVItem{September 2016 and September 2017, \textit{Teachers Assistant}, BK1 (Basic course 1 in data processing and statistics) }{I }

%------------------------------------------------

\CVItem{October 2015 - December 2018, \textit{Science Show Representative}}{One of the founding students of RUC's science show. Responsible for preparation and planning of new experiments and shows at various Danish high schools. The competences that i have gained through being a science show representative are the following. Comfortable speak in public and being able to engage audiences. I have didactic competence in setting the explanation of concepts in science for different age groups. I am able to uphold agreements and meet deadlines. I work well in a team environment.}

%------------------------------------------------


%\Sep % Extra whitespace after the end of a major section

%----------------------------------------------------------------------------------------
%	COMMUNICATION SKILLS
%----------------------------------------------------------------------------------------

\CVItem{June 2017 - June 2018, \textit{Co-creator at Roskilde Festival}, Volunteering}{
I have helped build and take down art exhibitions, stages and stage areas, working 100+ hours over a 14 day period. My jobs have also included driving telehandlers transporting heavy goods and painting various installations.
}

%\Sep
%\clearpage

%\userinformation % Print your information in the left column

%\framebreak % End of the first column

%------------------------------------------------

%\CVItem{2014, \textit{Poster}, Annual Business Conference (Oregon)}{As part of the course work for BUS320, I created a poster analyzing several local businesses and presented this at a conference.}

%------------------------------------------------

\Sep

\clearpage % Start a new page

\userinformation % Print your information in the left column

\framebreak % End of the first column

%----------------------------------------------------------------------------------------
%	SKILLS
%----------------------------------------------------------------------------------------

\CVSection{Technical skills}

%------------------------------------------------

\CVItem{Programming and systems}
{\begin{tabular}{p{0.2\textwidth} p{0.2\textwidth} p{0.2\textwidth}}
 \bluebullet GNU Octave &  \bluebullet MatLab\\
\end{tabular}}


%------------------------------------------------

\CVItem{Computer Software}
{\begin{tabular}{p{0.2\textwidth} p{0.2\textwidth} p{0.2\textwidth}}
 \bluebullet Linux &  \bluebullet LateX & \bluebullet Microsoft Office\\
\end{tabular}}

\Sep

\CVSection{Language}
{\begin{tabular}{p{0.2\textwidth} p{0.2\textwidth} p{0.2\textwidth}}
\CVItem{Danish}{ - Professional} & \CVItem{English}{ - Mothertounge}  &
\end{tabular}}

\Sep

%\CVSection{References}
%\CVItem{ (Master Thesis and Project Supervisor)}{Professor ()\\
%Dept. of Science and Environment
%Roskilde University,\\
%Denmark \\
%\emph{Contact info can be provided if needed}}
%------------------------------------------------

\Sep % Extra whitespace after the end of a major section

%----------------------------------------------------------------------------------------
%	NEW PAGE DELIMITER
%	Place this block wherever you would like the content of your CV to go onto the next page
%----------------------------------------------------------------------------------------

%\clearpage % Start a new page

%\userinformation % Print your information in the left column

%\framebreak % End of the first column

%----------------------------------------------------------------------------------------
%	AWARDS
%-----------------------------------------------------------------------------

%\CVSection{Awards}

%----------------------------------------------------------------------------------------
%	INTERESTS
%----------------------------------------------------------------------------------------

\CVSection{Interests}

%------------------------------------------------

\CVItem{Professional}{Mathematical Modeling, Computional Physics and programming, Experimental Physics: Small angle x-ray scattering, Grazing incidence small angle scattering, Dynamic light scattering, Reflectometry}

%------------------------------------------------

\CVItem{Personal}{I am a 29 year old Australian who is both an Australian and Danish citizen and i have lived in Danmark for 10 years. My hobbies are running, bikepacking and playing bass guitar. A big part of my summer is volunteering at Roskilde festival building the Pavilion stage and the Pavilion area. When winter rolls around i enjoy skiing.}


%------------------------------------------------

\Sep % Extra whitespace after the end of a major section

%----------------------------------------------------------------------------------------

\end{document}
